%!TEX root = ./../../springer-mono-format/book.tex
%!TEX root = ../../thesis-trier.tex

Here, we investigate how 
applying operations to 
constraint languages
affects the computational complexities of resulting constrained problems. Specifically, we investigate the following operations: 
the prefix/suffix/factor operations (\SEC~\ref{constr_sync:subsec:factor}),
set union (\cref{constr_sync:subsec:union}),
intersection with a regular
ideal language (Section~\ref{constr_sync:subsec:ideal}), applying morphism (Section~\ref{constr_sync:subsec:morphisms}),
and concatenating specific languages to an arbitrary constraint language (Section~\ref{constr_sync:subsec:concat}).
The operations from \cref{constr_sync:subsec:factor} 
preserve the complexity,
those from \cref{constr_sync:subsec:union,constr_sync:subsec:ideal,constr_sync:subsec:morphisms}
mostly decrease the complexity,
and those from \cref{constr_sync:subsec:concat}
increase the complexity.

\subsection{Factor Sets and Constraint Languages}
\label{constr_sync:subsec:factor}

%!TEX root = ./../../../thesis-trier.tex

Our first result is a consequence of Theorem~\ref{prelim:thm:sync_words_is_ideal}. Although it has a simple proof,
it is quite important for many of the following results.

\begin{theorem}
	\label{constr_sync:thm:L_in_FactLbar}
	Let $L, L' \subseteq \Sigma^*$ with $L \subseteq \factor{L'}$.
	If $\semiautomat{A}{Q}{\Sigma}{\delta}$ 
	has a synchronizing word in $L$, then it has one in $L'$.
\end{theorem}
\begin{proof}
	Suppose $u \in L$ synchronizes $\mathcal A$.
	By assumption there exist $v, w \in \Sigma^*$
	such that $vuw \in L'$.
	By Theorem~\ref{prelim:thm:sync_words_is_ideal} the
	latter word synchronizes $\mathcal A$ as well.
\end{proof}

Note that $L \subseteq \factor{L'}$ is equivalent to $\factor{L} \subseteq \factor{L'}$.
%As $L \subseteq L'$ implies $\factor{L} \subseteq \factor{L'}$
%and $\factor{\factor{L}} = \factor{L}$, the next is implied.
This gives the following consequence of Theorem~\ref{constr_sync:thm:L_in_FactLbar}.

\begin{corollary}
	\label{constr_sync:cor:factors_sets_included}
	Let $L, L' \subseteq \Sigma^*$ with $\factor{L} = \factor{L'}$.
	Then every automaton $\mathcal A$ has a synchronizing word in $L$
	if and only if it has one in~$L'$.
\end{corollary}

For any $L \subseteq \Sigma^*$ we have
$\factor{L} = \factor{\pref{L}} = \factor{\suff{L}}$. So with Corollary~\ref{constr_sync:cor:factors_sets_included},
we conclude that a given automaton has a synchronizing word in one of
the languages $L$, $\factor{L}$, $\pref{L}$, or $\suff{L}$ if and only if it has one in all of these languages.
As a consequence, we obtain that the corresponding constrained synchronization problems
are equivalent for $\leACZEROm$-reductions (in fact,
the input can be passed through without any modifications
and so these problems are also equivalent
for reductions by projections as considered in~\cite{valiant1982rfducibility}).
%, but $\leACZEROm$-reductions is sufficient for our applications). 

\begin{corollary}
	\label{gen_results:cor:factors_etc_complexity}
	Let $L \subseteq \Sigma^*$ be a regular language. 
	The following problems are equivalent for~$\leACZEROm$-reductions.
	\begin{enumerate}
		\item $\constrsync{L}$
		\item $\constrsync{\factor{L}}$
		\item $\constrsync{\pref{L}}$
		\item $\constrsync{\suff{L}}$
	\end{enumerate}
	Furthermore, if for two languages $L, L'\subseteq \Sigma^*$ the factor, prefix, or suffix sets %closures
	are equal, then $\constrsync{L} \equiv_m^{\AC^0} \constrsync{L'}$.
\end{corollary}

\begin{corollary} \label{constrsync:genresults:cor:union-fact}
	Let $L, L' \subseteq \Sigma^*$ be regular languages.
	Then we have  $\constrsync{L \cup L'} \eqACZEROm \constrsync{\fact{L}\cup L'}$
\end{corollary}
\begin{proof}
	This follows from \cref{gen_results:cor:factors_etc_complexity}
	since $\fact{L \cup L'} = \fact{\fact{L}\cup L'}$.
\end{proof}

\begin{corollary} \label{constrsync:genresults:cor:left-right-ideal}
	Let $L \subseteq \Sigma^*$ be a regular language
	containing a non-empty left or right ideal language.
	Then we have $\constrsync{L} \eqACZEROm \constrsync{\Sigma^*}$.	
\end{corollary}
\begin{proof}
	Without loss of generality, 
	we have $I \subseteq L$
	for some non-empty left ideal language~$I \subseteq \Sigma^*$ (which might be non-regular).
	Pick some $u \in I$ and consider a word $v \in \Sigma^*$.
	Then $vu \in I$
	so that $v \in \fact{I}$.
	This gives $\Sigma^* \subseteq \fact{I} \subseteq \fact{L}$.
	Thus $\fact{L} = \Sigma^*$
	and the claim follows from \cref{gen_results:cor:factors_etc_complexity}.
\end{proof}

\begin{corollary}
	Let $I, J \subseteq \Sigma^*$ be a left and a right ideal, respectively,
	which are both non-empty.
	Then $\constrsync{I\cap J} \eqACZEROm \constrsync{\Sigma^*}$.
\end{corollary}
\begin{proof}
	By assumption, we have $u \in I$ and $v \in J$ for some $u,v \in \Sigma^*$.
	Let $w \in \Sigma^*$.
	Then we have $vwu \in I$ since~$I$ is a left ideal 
	and $vwu \in J$ since~$J$ is a right ideal.
	Thus $vwu \in I \cap J$.
	This gives $\Sigma^* \subseteq \fact{I \cap J}$
	and the claim follows from \cref{gen_results:cor:factors_etc_complexity}. 
\end{proof}

\subsection{The Union of Constraint Languages}
\label{constr_sync:subsec:union}

%!TEX root = ../../../thesis-trier.tex

Here, in~\cref{constr_sync:thm:union}, we state that certain complexity classes
are preserved when considering the union
of two constraint languages. \cref{constrsync:genresults:cor:subset}
gives further conditions
for complexity preservation w.r.t.~union
and concatenation.
% when the complexity is preserved 
%for unions and concatenation of languages.

\begin{theorem}
	\label{constr_sync:thm:union}
	Let $\mathcal C$ be one of the complexity classes $\AC^0$, $\NC^1$, \LOGSPACE, \NL, \PTIME or \NP.
	Suppose $\constrsync{L}$ and $\constrsync{L'}$ are both in $\mathcal C$ for regular languages $L,L' \subseteq \Sigma^*$.
	Then  the problem $\constrsync{L \cup L'}$ is in $C$.
\end{theorem}
\begin{proof}
%Given two computing mechanism (circuits or Turing machines)
%$M_1$ and $M_2$ for \constrsync
For \LOGSPACE, \NL, \PTIME, and \NP\ suppose
$M_1$ and $M_2$ are Turing machines deciding $\constrsync{L}$
and $\constrsync{L'}$, respectively, for an input automaton
by halting and writing~$1$ for a \rYES-instance
or~$0$ for a \rNO-instance.
Let $M$ be a Turing machine that, for an input automaton $\mathcal A$,
first simulates $M_1$ with input~$\mathcal A$ yielding output~$r_1$,
then $M_2$ with input $\mathcal A$ yielding output $r_2$, and finally outputs $\max\{r_1, r_2\}$,
which equals one if and only if $\mathcal A$ has a synchronizing word in $L \cup L'$.  
We can combine the results $r_1$ and $r_2$  (repositioning reading head, erasing the working tape while storing the output of $M_1$, compute
$\max\{r_1, r_2\}$) 
by operating deterministically and using only constant space;
and if~$M_1$ and~$M_2$ obey the convention to
position the working tape head next to the output, erase
everything else and  move the reading head back to the start,
the overhead
needed to simulate both machines and combine  their results is actually constant in time and space.
Hence, if $M_1$ and $M_2$ both work in deterministic logspace, nondeterministic logspace,
deterministic polynomial time, or nondeterministic polynomial time,
also $M$ works in deterministic logspace, nondeterministic logspace,
deterministic polynomial time, or nondeterministic polynomial time, respectively.

For $\AC^0$ or $\NC^1$, we can combine the outputs
of two circuits from respective circuit families in $\AC^0$ or $\NC^1$ for $\constrsync{L}$
and $\constrsync{L'}$ by an OR-gate
to obtain an \LOGSPACE-uniform circuit family in $\AC^0$ or $\NC^1$, respectively, 
for $\constrsync{L \cup L'}$.
\end{proof}

Next, we state a condition such that the complexity is preserved when taking the union.

\begin{corollary} \label{constrsync:genresults:cor:subset}
	Let $L, L' \subseteq \Sigma^*$ be regular languages
	such that $L' \subseteq \fact{L}$.
	Then we have $\constrsync{L} \eqACZEROm \constrsync{L\cup L'}$.
\end{corollary}
\begin{proof}
	By \cref{constrsync:genresults:cor:union-fact}
	we have $\constrsync{L\cup L'} \eqACZEROm \constrsync{\fact{L} \cup L'}$
	and the assumption gives $\fact{L} \cup L' = \fact{L}$.
	By  \cref{gen_results:cor:factors_etc_complexity}
	we have $\constrsync{\fact{L} \cup L'} \eqACZEROm \constrsync{L}$ (the claim
	also follows directly from \cref{gen_results:cor:factors_etc_complexity}
	by noting that $\fact{L \cup L'} = \fact{L} \cup \fact{L'} = \fact{L}$).
\end{proof}

By \cref{constrsync:genresults:cor:subset}, for a right or left ideal~$I$ that is regular
and any regular $L \subseteq \Sigma^*$
we have $\constrsync{I \cup L} \eqACZEROm \constrsync{\Sigma^*}$
since $\factor{I} = \Sigma^*$.

\begin{remark}
	Let $L = \{a,b\}^* c \{a,b\}^*$ for distinct $a,b,c\in \Sigma$. 
	By \crefitem{constr_sync:thm:introductory_examples}{constrsync:ex:item:abcL}, $\constrsync{L}$
	is \PSPACE-complete. We have $\Sigma^* \setminus L = \{ u \in \Sigma^* : |u|_c = 0 \text{ or } |u|_c \ge 0 \}$
	and $\fact{\Sigma^* \setminus L} = \{a,b,c\}^*$.
	By \crefitem{constr_sync:thm:introductory_examples}{constrsync:ex:item:ab}, the problem
	$\constrsync{\Sigma^* \setminus L}$ is \NL-complete.
	Thus, in general, the computational complexity of the constrained synchronizing problem
	is not preserved when taking the complement.
\end{remark} 

\subsection{The Ideal Intersection Theorem}
\label{constr_sync:subsec:ideal}

%!TEX root = ./../../../thesis-trier.tex

Here, we show the ideal intersection theorem (\cref{constr_sync:gen_results:thm:ideal_separation}).
Later, we mostly use this result
to deduce \PSPACE-hardness from known \PSPACE-hard problems, especially
for constraint languages written as unions (see Remark~\ref{constrsync:compl:ideal_intersection_for_unions} for details).

\begin{theorem}[Ideal Intersection Theorem]
	\label{constr_sync:gen_results:thm:ideal_separation}
	Let $I \subseteq \Sigma^*$ be a fixed ideal language that is regular.
	Suppose $L \subseteq \Sigma^*$ is any regular language,
	then
	$$
	\constrsync{L \cap I} \leACZEROm \constrsync{L}.
	$$
\end{theorem}
\begin{proof}
	By \cref{prelim:thm:sync_words_is_ideal}
	there exists an automaton $\mathcal A_I$
	such that $I = \Syn(\mathcal A_I)$.
	Let $\semiautomat{A}{Q}{\Sigma}{\delta}$ be an input to $\constrsync{L\cap I}$.
	By~\cref{prelim:eqn:product_aut_sync_words}, 
	it is $\Syn(\mathcal A) \cap I = \Syn(\mathcal A \times \mathcal A_I)$,
	which gives	
	\begin{equation}\label{constrsync:eqn:ideal-int}
	\Syn(\mathcal A) \cap (L \cap I) \ne \emptyset
	\Leftrightarrow 
    \Syn(\mathcal A \times \mathcal A_I) \cap L \ne \emptyset,
	\end{equation}
	i.e., $\mathcal A \times \mathcal A_I$ has a synchronizing
	word in $L$ if and only if $\mathcal A$ has a synchronizing word in $L \cap I$.
	As $I$ is fixed,
	we can construct the product automaton $\mathcal A \times \mathcal A_I$
	%for a fixed automaton $\mathcal A_I$ 
	by a circuit family in $\AC^0$ 
	 as we only have to add transitions
	and states to~$\mathcal A$ determined by a fixed automaton $\mathcal A_I$.
	So, by~\cref{constrsync:eqn:ideal-int} 
	constructing $\mathcal A \times \mathcal A_I$ gives a valid reduction to \constrsync{L}.	
\end{proof}

\begin{example}
	Let $\Sigma = \{a,b,c\}$ and consider the ideal language $I = \Sigma^* c^{n} \Sigma^*$
	that is regular.
	Let $L_1 = \{ u \in \Sigma^* : |u|_c \le n \}$,
	$L_2 = \{ u \in \Sigma^* : |u|_c = n \}$,
	and $L_3 = \{ u \in \Sigma^* : m \le |u|_c \le n \}$
	for $n,m \ge 1$ with $m \le n$.
	Then we have
	$L_1 \cap I = L_2 \cap I = L_3 \cap I = \{a,b\}^* c^n \{a,b\}^*$.
	From \crefitem{constr_sync:thm:introductory_examples}{constrsync:ex:item:abcL} 
	and \cref{constr_sync:gen_results:thm:ideal_separation} 
	we get that $\constrsync{L_i}$
	for $i\in \{1,2,3\}$
	are all \PSPACE-complete problems.
\end{example}

\begin{corollary}
	Let $L \subseteq \Sigma^*$ be regular
	and $u_1, \ldots, u_n \in \Sigma$.
	Then we have 
	\begin{align*}
	\constrsync{L\cap \Sigma^*u_1\cdots u_n\Sigma^*} & \leACZEROm \constrsync{L}, \\
	\constrsync{L\cap \Sigma^*u_1 \Sigma^* \cdots \Sigma^* u_n \Sigma^*} & \leACZEROm \constrsync{L}.
	\end{align*}
\end{corollary}
\begin{proof}
	The languages
	 $\Sigma^*u_1\cdots u_n\Sigma^*$
	 and $\Sigma^*u_1 \Sigma^* \cdots \Sigma^* u_n \Sigma^*$
	 are regular and ideals.
	 So, the claim follows from \cref{constr_sync:gen_results:thm:ideal_separation}.		
\end{proof}

\IN{principal ideal}
Mostly, we apply Theorem~\ref{constr_sync:gen_results:thm:ideal_separation}
with \emph{principal ideals} $\Sigma^* u \Sigma^*$
for $u \in \Sigma^*$.

As pointed out in the next remark, for a union the problem either has the same complexity as each part, or its complexity is bounded from below by the intersection of one part with a principal ideal.

\begin{remark}\label{constrsync:compl:ideal_intersection_for_unions}
	Consider two regular languages $L_1, L_2 \subseteq \Sigma^*$.
	There exists $u \in L_2$
	such that $u \notin \fact{L_1}$
	or $L_2 \subseteq \fact{L_1}$.
	By \cref{constrsync:genresults:cor:subset},
	in the former case, 
	we have
	\[
	\constrsync{L_1} \eqACZEROm \constrsync{L_1\cup L_2}
	\]
	and in the latter case, since then
	\[
	( L_1 \cup L_2 )  \cap \Sigma^* u \Sigma^* = L_2 \cap \Sigma^* u \Sigma^*,
	\]
	from \cref{constr_sync:gen_results:thm:ideal_separation}
	we get 
	$\constrsync{L_2 \cap \Sigma^*u \Sigma^*} \leACZEROm \constrsync{L_1\cup L_2}$.
	%	This scheme is often use in hardness proofs.
	%	However, finding an appropriate word $u \in \fact{L_2}$ 
	%	such that a provable 
	%	hard problem $\constrsync{L_2 \cap \Sigma^*u \Sigma^*}$
	%	results is the tricky part.
\end{remark}

Because we often use \cref{constr_sync:gen_results:thm:ideal_separation} to ``separate
a hard part out'' of a constraint language we called it the Ideal Separation Theorem
in~\cite{HoffmannCocoonExtended}, but we changed the name here because it involves an intersection
and we thought the new name to be more appropriate.

Lastly, we show how the Ideal Intersection Theorem  can be of help
in a heuristic way by employing it to exclude problematic
cases from a construction. This usage might not always be possible (in fact, in most cases it is not 
possible),
but we give the following example to also build some intuition
on various ways in which the Ideal Intersection Theorem 
can be applied (more examples of its application 
will follow in the rest of this work).

\begin{remark} \label{constrsync:rem:construction-ideal-intersection}
	Consider $L = \{bb,ba\}^*a\Gamma^*$ for $|\Gamma|\le 1$.
	We use a reduction from \SIS
	to show \PSPACE-hardness of $\constrsync{L}$.
	However, the construction we give
	does not work straightforward way for~$L$,
	but we then argue that it can be
	used for a modified problem obtained by intersection with an appropriate ideal.
	We then obtain \PSPACE-hardness of $\constrsync{L}$
	from \cref{constr_sync:gen_results:thm:ideal_separation}.
	
	Let $\semiautomat{A}{Q}{\{a,b\}}{\delta}$ and $S \subseteq Q$.
	Set $\semiautomat{A'}{Q \times \{\varepsilon,b\} \cup \{t\}}{\{a,b\}}{\delta'}$
	where~$\delta'(t, x) = t$ and
	\[
	\delta'((q, u), x) =
	\begin{cases}
		t                           & \text{if } u = \varepsilon \text{ and } x = a, \\
		(q, b) & \text{if } u = \varepsilon \text{ and } x = b, \\
		(\delta(q, a), \varepsilon) & \text{if } ux = bb, \\
		(\delta(q, b), \varepsilon) & \text{if } ux = ba, \\
		(q, \varepsilon)            & \text{if } ux \notin \{bb,ba\} \text{ and } q \notin S, \\
		t							& \text{if } ux \notin \{bb,ba\} \text{ and } q \in S, 		
	\end{cases}
	\]
	for $u \in \{\varepsilon, b\}$, $q \in Q$ and $x \in \{a,b\}$.
	
	First, suppose $u_1 \cdots u_n a v$
	with $u_1, \ldots, u_n \in \{bb,ba\}$ and $v \in \Gamma^*$ 
	synchronizes~$\mathcal A'$.
	Then it must map all states to~$t$.
	Consider the state $\delta'((q, \varepsilon), u_1\cdots u_n)$
	for some $q \in Q$.
	Then it must have the form $(p, \varepsilon)$
	for some $p \in Q$.
	Assume $p \notin S$.
	Then $\delta'((p, \varepsilon), a) = (p, \varepsilon)$
	and since $|\Gamma|\le 1$
	the word~$v$ cannot map $(p, \varepsilon)$ to~$t$.
	Thus we must have $p \in S$
	and then the word $v_1 \cdots v_n$
	with $v_i = a$ if $u_i = bb$
	and $v_i = b$ if $u_i = ba$
	maps~$Q$ into~$S$.	
	
	If $u_1 \cdots u_n$ for $u_1,\ldots,u_n \in \{a,b\}$
	maps~$Q$ into~$S$,
	then $v_1 \dots v_n a$ with $v_i = bb$ if $u_i = a$
	and $v_i = ba$ if $u_i = b$ for all $i \in \{1,\ldots,n\}$
	maps all states of~$Q \times \{\varepsilon\}$ to~$t$.
	However, if $u_1 = \cdots = u_n = bb$,
	then $\delta'(Q \times \{b\}, u_1\cdots u_n) \subseteq Q \times \{b\}$
	and thus $\delta'(Q \times \{b\}, u_1 \cdots u_n a) \cap \delta(Q, a) \times \{\varepsilon\} \ne \emptyset$.
	Now no word from~$\Gamma^*$ can map this state to~$t$.
	But this problem does not arise if we read~$ba$
	at least once, since this word maps all states in $Q \times \{b\}$
	to~$t$.
	Now, note that $L \cap \Sigma^* baa \Sigma^* = \{bb,ba\}baa\Gamma^*$.
	Thus, by \cref{constr_sync:gen_results:thm:ideal_separation} we obtain \PSPACE-hardness for $\constrsync{L}$
	if we can show \PSPACE-hardness
	for $\constrsync{\{bb,ba\}baa\Gamma^*}$
	and to establish \PSPACE-hardness of this problem
	we can use the above construction since the problematic case discussed above
	does not arise for $\{bb,ba\}baa\Gamma^*$.
\end{remark}

%\subsection{Homomorphic Mappings}
\subsection{Applying Morphisms and Solid Codes}
\label{constr_sync:subsec:morphisms}

%!TEX root = ../../../thesis-trier.tex

Here, in \cref{thm:hom_lower_bound_complexity},
we show that transforming a constraint language
yields a problem whose complexity is
at most that of the original problem w.r.t.~$\leACZEROm$-reductions.
In \cref{thm:constr_sync_hom_strongly_self_sync}
we show that for coding morphisms of solid codes
the complexity is preserved.

\begin{theorem}
	\label{thm:hom_lower_bound_complexity}
	Let $\varphi \colon \Sigma^* \to \Gamma^*$ be a fixed morphism
	and $L \subseteq \Sigma^*$.
	Then \[ \constrsync{\varphi(L)} \le_m^{\AC^0} \constrsync{L}. \]
\end{theorem}
\begin{proof}
	For the $\leACZEROm$-reduction,
	let $\semiautomat{A}{Q}{\Gamma}{\delta}$ be an input automaton. 
	We want to know if $\mathcal A$ is synchronizing with respect to~$\varphi(L)$.
	
	\medskip
	\noindent\underline{Construction:} Let $\semiautomat{A'}{Q}{\Sigma}{\delta'}$
	with $\applytrans{\delta'}{q}{x} = \applytrans{\delta}{q}{\applyfunc{x}{\varphi}}$ for $q \in Q$ and $x \in \Sigma$.
	We obtain a representation of $\mathcal A'$
	by transition matrices by multiplying
	for each $x \in \Sigma$
	the respective transition matrices
	of the letters in~$\applyfunc{x}{\varphi}$,
	which is doable by a \LOGSPACE-uniform circuit family in $\AC^0$ as $\varphi$ is fixed and matrix multiplication for a fixed number of matrices is in $\AC^0$ by \cref{prelim:circuit:ex:AC0mtx}.
		
	\medskip 
	\noindent\underline{Verification:} We have
	\begin{align*}
		                  \exists u \in \varphi(L) \colon |\delta(Q, u)| = 1 
		& \Leftrightarrow \exists v \in L \colon |\delta(Q, \varphi(v))| = 1 \tag{$\varphi(v) = u$} \\
		& \Leftrightarrow \exists v \in L \colon |\delta'(Q, v)| = 1  \tag{$\delta'(Q, v) = \delta(Q, \varphi(v))$}.  
	\end{align*}
So $\mathcal A'$ has a synchronizing word in $L$ if and only if
$\mathcal A$ has a synchronizing word in $\applyfunc{L}{\varphi}$.
\end{proof}

\begin{example}
	Let $L = \{a,b\}^* \cup \{a,c\}^* \cup \{b,c\}^*$.
	Then $\pi_{\{a,b\}}(L) = \{a,b\}^*$,
	so that $\constrsync{L}$ is \NL-hard
	with \cref{thm:hom_lower_bound_complexity}
	and \crefitem{constr_sync:thm:introductory_examples}{constrsync:ex:item:ab}.
	With \cref{constr_sync:thm:union} 
	we also get membership in \NL.
\end{example}	

For coding morphisms of solid codes,
we have a stronger statement than 
Theorem~\ref{thm:hom_lower_bound_complexity}.

\begin{theorem}
	\label{thm:constr_sync_hom_strongly_self_sync}
	Let $\varphi \colon \Sigma^* \to \Gamma^*$
	be a coding morphism such that\footnote{Recall that $\Sigma$ is finite, and so $C$ is finite. Also coding morphism
		are injective by definition.}
	$C = \varphi(\Sigma)$ 
	is a solid code.
	Then, for each regular $L \subseteq \Sigma^*$ we have 
	$
	\constrsync{L} \eqACZEROm \constrsync{\varphi(L)}.
	$
\end{theorem}
\begin{proof}
	By Theorem~\ref{thm:hom_lower_bound_complexity}, we have 
	$\constrsync{\varphi(L)} \leACZEROm \constrsync{L}$.
	Next, we show $\constrsync{L} \leACZEROm \constrsync{\applyfunc{L}{\varphi}}$
	by utilizing the Levenshtein-Romanov mappping from \cref{constr_sync:subsec:levenshtein-romanov-map}
    w.r.t.~the language~$C$.
	Now, let $\semiautomat{A}{Q}{\Sigma}{\delta}$ be an input automaton.
	Suppose $\Sigma = \{a_1, \ldots, a_k\}$
	and set $u_i = \applyfunc{a_i}{\varphi}$ for $i \in \{1,\ldots,k\}$.
	Then $C = \{u_1, \ldots, u_k\}$.	
	
	\begin{figure}[htbp]
		\centering		
		\begin{subfigure}{0.33\textwidth}
			\centering
	\begin{tikzpicture}[>=stealth,shorten >=1pt,auto]
	\node[state] (q0) at (90:2.5cm) {$q_0$};
	\node[state] (q1) at (210:2.5cm) {$q_1$};
	\node[state] (q2) at (330:2.5cm) {$q_2$};
	
	\path[->]
	(q0) edge[bend left=10] node[right] {$a_1$} (q2)
	(q2) edge[bend left=10] node[below] {$a_1$} (q1)
	(q1) edge[bend left=5]  node[left] {$a_1$}  (q0);
	
	\path[->] (q1) edge[bend left=10] node {$a_2$} (q2)
	(q2) edge[bend left=10] node {$a_2$} (q0)
	(q0) edge [loop above] node {$a_2$} (q0);
\end{tikzpicture}
			\caption{A base automaton~$\mathcal A$}
		\end{subfigure}
		\hfill
		\begin{subfigure}{0.55\textwidth}
		\centering
	\begin{tikzpicture}[>=stealth,shorten >=1pt, node distance=0.97cm, auto]
	
	%\tikzset{state/.style={}}
	\tikzset{state/.style={inner sep=2pt}}
	
	%%%%%%%%%%%%%%%%%%%%%%%%%%%%%%%%%%%%
	% TOP LAYER (epsilon)
	%%%%%%%%%%%%%%%%%%%%%%%%%%%%%%%%%%%%
	\node[state] (q0e) at (90:2.7cm) {$(q_0,\varepsilon)$};
	\node[state] (q1e) at (210:2.7cm) {$(q_1,\varepsilon)$};
	\node[state] (q2e) at (330:2.7cm) {$(q_2,\varepsilon)$};
	
	%%%%%%%%%%%%%%%%%%%%%%%%%%%%%%%%%%%%
	% MIDDLE LAYER (a)
	%%%%%%%%%%%%%%%%%%%%%%%%%%%%%%%%%%%%
	\node[state] (q0a) [below=of q0e] {$(q_0,a)$};
	\node[state] (q1a) [below=of q1e] {$(q_1,a)$};
	\node[state] (q2a) [below=of q2e] {$(q_2,a)$};
	
	%%%%%%%%%%%%%%%%%%%%%%%%%%%%%%%%%%%%
	% LOWER LAYER (aa)
	%%%%%%%%%%%%%%%%%%%%%%%%%%%%%%%%%%%%
	\node[state] (q0aa) [below=of q0a] {$(q_0,aa)$};
	\node[state] (q1aa) [below=of q1a] {$(q_1,aa)$};
	\node[state] (q2aa) [below=of q2a] {$(q_2,aa)$};
	
	\path[->]
	(q0e) edge (q0a)
	(q0a) edge (q0aa)
	(q0aa) edge[in=230,out=260, looseness=9] (q0aa)
	(q0aa) edge[dotted, in=280, out=310, looseness=9] (q0aa)
	
	(q1e) edge  (q1a)
	(q1a) edge  (q1aa)
	(q1aa) edge[loop below] ()
	
	(q2e) edge  (q2a)
	(q2a) edge  (q2aa)
	(q2aa) edge[loop below] ();		
	
	\path[->]
	(q0e)  edge[dashed, in=80, out=50, looseness=9] (q0e)
	(q0e)  edge[dotted, in=130, out=100, looseness=9] (q0e)
	(q0a)  edge[dotted, bend left=20] (q0e)
	(q0aa) edge[dashed, bend right=10] (q2e);
	
	\path[->]
	(q1e)  edge[dashed, in=80, out=50, looseness=9] (q1e)
	(q1e)  edge[dotted, in=130, out=100, looseness=9] (q1e)
	(q1a)  edge[dotted, bend left=20] (q1e)
	(q1aa) edge[dashed, bend left=65] (q0e);
	
	\path[->]
	(q2e)  edge[dashed, in=80, out=50, looseness=9] (q2e)
	(q2e)  edge[dotted, in=130, out=100, looseness=9] (q2e)
	(q2a)  edge[dotted, bend right=20] (q2e)
	(q2aa) edge[dashed, bend left=30] (q1e);		
	
	\path[->] 
	(q0a) edge[dashed, bend right=10] (q2e)
	(q2a) edge[dashed,bend left=30] (q1e)
	(q1a) edge[dashed,bend left=55] (q0e);	
	
	\path[->]
	(q1aa) edge[dotted, bend right=10] (q2e)
	(q2aa) edge[dotted, bend right=50] (q0e);		
\end{tikzpicture}
			\caption{Constructed automaton as in the proof.}
		\end{subfigure}		
		\caption[Construction from proof of \cref{thm:constr_sync_hom_strongly_self_sync}.]{Construction from the proof of \cref{thm:constr_sync_hom_strongly_self_sync}
		    for the solid code $\{ab, aac\}$
			and base automaton shown on the left side. Solid, dashed, and dotted lines represent transitions
			labeled with $a$, $b$, and~$c$, respectively.}
	\label{constrsync:fig:construction-solid-code}
	\end{figure}	
	
	\medskip
	\noindent\underline{Construction:} Define $\semiautomat{A'}{Q \times (\pref{C} \setminus C)}{\Gamma}{\delta'}$
	where, for $(q, u) \in Q \times \pref{C} \setminus C$
	and $x \in \Sigma$,
	\begin{equation}\label{constr_sync:eqn:aut_solid_code}
		\applytrans{\delta'}{(q,u)}{x} = \begin{cases}
			(q, \sigma_C(ux)) & \text{if } \sigma_C(ux) \notin C, \\
			(\applytrans{\delta}{q}{a_i}, \varepsilon) & \text{if } \exists i \in \{1,\ldots,k\} : \sigma_C(ux) = u_i. 
		\end{cases}
	\end{equation}
	By encoding the proper prefixes in the second component
	we keep track of the prefix scanned in the input so far.	
	See \cref{constrsync:fig:construction-solid-code} for an example construction
	for the solid code $C = \{ab,aac\}$.
	Note that in the example $\delta'((q, a), b) = \delta'((q,aa), b) = (\delta(q, a_1), \varepsilon)$
	since $\sigma_C(ab) = \sigma_C(aab) = ab$.
		
	
	\noindent\underline{Properties:} We give the proofs for Properties~\ref{constr_sync:solid_code_morphism:apply_codeword_proper_prefix} and~\ref{constr_sync:solid_code_morphism:apply_codeword},
	Property~\ref{constr_sync:solid_code_morphism:wordsvarphi} follows from Property~\ref{constr_sync:solid_code_morphism:apply_codeword} by induction
	as $\applyfunc{v}{\varphi} \in C^+$ for $v \in \Sigma^+$.
	\begin{enumerate}
		\item \label{constr_sync:solid_code_morphism:apply_codeword_proper_prefix}
		For $(q, u) \in Q \times ( \pref{C} \setminus C )$ and $v \in \pref{C} \setminus C$
		we have $\sigma_C(uv) \in \pref{C} \setminus C$ and $\applytrans{\delta'}{(q,u)}{v}  = (q, \sigma_C(uv))$.		
		
		\begin{proofitem}{\ref{constr_sync:solid_code_morphism:apply_codeword_proper_prefix}} 
		We show the claim by induction on $v$.
		As $\sigma_C(u) = u$, the claim is true for $v = \varepsilon$.
		Now, suppose $v = wx$ for $w \in \Sigma^*$ and $x \in \Sigma$.		
		As $C$ is overlap-free, 
		$wx \in \pref{C} \setminus C$ implies 
		$\sigma_C(uwx) \notin C$ and so $\sigma_C(uwx) \in \pref{C} \setminus C$;
		this implication is noted in \crefitem{constrsync:prop:sigma-solid-code-char}{constrsync:prop:sigma-solid-code-char:2}.
		By~\crefitem{constr_sync:prop:sigma_L_properties}{constr_sync:sigme_L:item:concatenation}
		we have $\sigma_C(\sigma_C(uw)x) = \sigma_C(uwx)$.
		So, for $q \in Q$, we obtain		
		\begin{align*}
			& \delta'((q,u), wx) \\
			& = \delta'(\delta'((q,u), w), x) & \tag{Def.~of $\delta'$} \\
			& = \delta'((q, \sigma_C(uw)), x) & \tag{Induction hypothesis} \\
			& = (q, \sigma_C(\sigma_C(uw)x)). & \tag{Eqn.~\eqref{constr_sync:eqn:aut_solid_code} as $\sigma_C(\sigma_C(uw)x) \notin C$} 
		\end{align*}
		Hence, we get $ \delta((q,u), wx) = (q, \sigma_C(uwx))$
        and the claim follows by induction.
     	\end{proofitem}
		
		\item \label{constr_sync:solid_code_morphism:apply_codeword}
		For $(q, u) \in Q \times ( \pref{C} \setminus C )$ and $i \in \{1,\ldots,k\}$ we have
		\[
		\applytrans{\delta'}{(q,u)}{u_i} = (\applytrans{\delta}{q}{a_i}, \varepsilon) 
		\]
		
		\begin{proofitem}{\ref{constr_sync:solid_code_morphism:apply_codeword}}
		Let $u_i = vx$ for $v \in \Sigma^*$ and $x \in \Sigma$.
		As $C$ is a prefix code, we have $v \in \pref{C} \setminus C$.		
		By \crefitem{constr_sync:prop:sigma_L_properties}{constr_sync:sigme_L:item:concatenation}
		we have $\sigma_C(\sigma_C(uv)x) = \sigma_C(uvx)$.
		By \crefitem{constrsync:prop:sigma-solid-code-char}{constrsync:prop:sigma-solid-code-char:3}
		as $vx \in C$, we find $\sigma_C(uvx) = vx = u_i$. Then
		\begin{align*}
			\delta'((q, u), vx) & = \delta'(\delta'((q,u), v), x)  & \tag{Def.~of $\delta'$} \\
			& = \delta'((q, \sigma_C(uv)), x) & \tag{Property~\ref{constr_sync:solid_code_morphism:apply_codeword_proper_prefix}} \\
			& = (\applytrans{\delta}{q}{a_i}, \varepsilon). & \tag{Eqn.~\eqref{constr_sync:eqn:aut_solid_code} as $\sigma_C(\sigma_C(uv)x)  = u_i$}
		\end{align*}
	    So, $\applytrans{\delta'}{(q,u)}{u_i} = (\applytrans{\delta}{q}{a_i}, \varepsilon)$.
		\end{proofitem}
		
		\item  \label{constr_sync:solid_code_morphism:wordsvarphi} For every $v \in \Sigma^+$ and $(q,u) \in Q \times (\pref{C} \setminus C)$,
		\begin{equation}\label{constr_sync:eqn:solid_code3}
			\applytrans{\delta'}{(q,u)}{\applyfunc{v}{\varphi}} = (\applytrans{\delta}{q}{v}, \varepsilon).
		\end{equation}
	\end{enumerate}	
	
	\medskip 
	\noindent\underline{Verification:} That the reduction is correct follows by the following
	equivalences:
	\begin{align*}
		& \exists w \in \varphi(L) \colon |\applytrans{\delta'}{Q \times (\pref{C}\setminus C)}{w}|  = 1 \\
		& \Leftrightarrow \exists v \in L \colon |\applytrans{\delta'}{Q \times (\pref{C}\setminus C)}{\varphi(v)}|  = 1 \\
		& \Leftrightarrow \exists v \in L \colon |\applytrans{\delta}{Q}{v} \times \{\varepsilon\}| = 1 \tag{Eqn.~\eqref{constr_sync:eqn:solid_code3}} \\
		& \Leftrightarrow \exists v \in L \colon |\applytrans{\delta}{Q}{v}| = 1. \qedhere
	\end{align*}
\end{proof}

\begin{remark} \label{constrsync:rem:alphabet_solid_special_case}
	The construction from \cref{constrsync:rem:alphabet} is a special case 
	of \cref{thm:constr_sync_hom_strongly_self_sync}
	as for $\Gamma, \Sigma$ with $\Gamma \subseteq \Sigma$
	the inclusion map $i \colon \Gamma \to \Sigma$
	given by $i(x) = x$ for $x \in \Gamma$
	is a coding morphism where $i(\Gamma) \subseteq \Sigma$ is a solid code.
\end{remark} 

\begin{remark}
	The construction from the proof of \cref{thm:constr_sync_hom_strongly_self_sync}
	with the transition function given by \cref{constr_sync:eqn:aut_solid_code}
	can be applied to any prefix code~$C$
	and yields a well-defined automaton~$\semiautomat{A'}{Q \times \pref{C}\setminus C}{\Sigma}{\delta'}$.	
	Can the assumption for~$C$ to be a solid code
	be weakened? 
	The crucial property for the reduction
	to work is that $\delta'(Q \times \pref{C} \setminus C, u_i) = \delta(Q, a_i) \times \{\varepsilon\}$
	for $i \in \{1,\ldots,k\}$ and $C = \{u_1, \ldots, u_k\}$.
	However, this is not guaranteed to hold if~$C$
	is merely a comma-free code.
	For example, consider $C = \{ u_1, u_2 \}$ with $u_1 = aab$ and $u_2 = bba$,
	which is a comma-free code, but not a solid code since $u_1$ and $u_2$ have a
	non-empty proper overlap.
	Then we have
	$
	\delta'(Q \times \{aa\}, u_2) = \delta(Q, a_1) \times \{a\}  	
	$
	and so $\delta(Q, a_1) \times \{a\} \subseteq \delta'(Q \times \pref{C} \setminus C, u_2)$.
	However, there are codes that are not solid codes for which it does hold.
	For example, for $C = \{ a_1, \ldots, a_{k-1}, a_k^m a_{k-1}\}$ for distinct $a_1, \ldots, a_k \in \Sigma$
	and $m \ge 1$.
\end{remark}

\begin{remark}
	Solid codes and the construction from the proof of  \cref{thm:constr_sync_hom_strongly_self_sync}
	can also be used to give an alternative proof of \cref{prelim:prop:sync_NL_complete}
	by first giving an automaton derived from the digraph
    by labeling the edges with letters and using an additional letter
    to realize the transitions going back to the start state.
    Then using a binary solid code
    and the construction from the proof of \cref{thm:constr_sync_hom_strongly_self_sync}
    we can reduce the alphabet to a binary alphabet.
    However, note that the words $bb,a,ba$ actually used in the proof of \cref{prelim:prop:sync_NL_complete}
    do not form a solid code.
\end{remark}

\begin{remark}
	We actually have $\Syn(\mathcal A') \cap \varphi(\Sigma^*) = \varphi(\Syn(\mathcal A))$
	in the construction from the proof of \cref{thm:constr_sync_hom_strongly_self_sync}.
\end{remark}

Coding morphism of solid codes are not the only morphisms
that preserve the complexity, as shown in the following remark.

\begin{remark}
	Consider the coding morphism $\varphi \colon \{a,b\}^* \to \{c,d, e\}^*$
	given by $\varphi(a) = cd$ and $\varphi(b) = ec$.
	Now $\varphi(\{a,b\}) = \{cd, ec\}$ is not a solid code.
	For $\semiautomat{A}{Q}{\{a,b\}}{\delta}$
	let 
	$\semiautomat{A'}{Q'}{\{c,d\}}{\delta'}$
	with $Q' = Q \times \{\varepsilon, c, e\}$
	where, for $q \in Q$, $u \in \{\varepsilon,c,d\}$, and $x \in \Sigma$,
	\[
	\delta'((q, u), x)
	= 
	\begin{cases}
		(\delta(q, a), \varepsilon)  & \text{if } u \in \{\varepsilon,c\} \text{ and } x = d, \\
		(\delta(q, b), \varepsilon)  & \text{if } u = e \text{ and } x = c, \\
		(q,x)                        & \text{if } u \in \{\varepsilon,c,e\} \text{ and } x \in \{c,e\} \text{ and } ux \ne ec \\
		%(q,e)                       & \text{if } u \in \{\varepsilon,e\} \text{ and } x = e,
		(q,u)                        & \text{if } u = e \text{ and } x = d.
	\end{cases}
	\]
	Then we have $\delta'((q, u), \varphi(v)) = (\delta(q, v), \varepsilon)$
	for $v \in \{a,b\}^+$ with $u \in \{\varepsilon, c, e\}$
	and $q \in Q$.
	This yields 
	\begin{align*}
		\delta'(Q', \varphi(a)) = \delta(Q, a) \times \{\varepsilon\}, \\
		\delta'(Q', \varphi(b)) = \delta(Q, b) \times \{\varepsilon\},	
	\end{align*}
	and	that for $|Q| \ge 2$ and regular $L \subseteq \{a,b\}^*$,
	the automaton~$\mathcal A$ has a synchronizing word in~$L$
	if and only if~$\mathcal A'$ has a synchronizing word in~$\varphi(L)$.
\end{remark}

\subsection{Concatenation and Complexity Lower Bounds}
\label{constr_sync:subsec:concat}

%!TEX root = ./../../../springer-mono-format/book.tex
%!TEX root = ./../../../thesis-trier.tex

In general, the computational complexity might increase or decrease
when concatenating languages.
For example, 
by \cref{constr_sync:thm:introductory_examples},
\cref{gen_results:cor:factors_etc_complexity}
and as $\fact{(a+b+c)^*c(a+b)^*} = \fact{(a+b)^*}$,
$\constrsync{(a+b)^*}$, $\constrsync{c(a+b)^*}$,
and $\constrsync{(a+b+c)^*c(a+b)^*}$ are 
complete for \NL, \PSPACE, and \NL, respectively.
Here, we identify conditions on concatenating languages
to obtain complexity lower bounds.
First, we give two results
for concatenating a language
to the front of a constraint language such that hardness is preserved, where
\cref{constrsync:prop:UanL} is related to languages contained in a subalphabet
and \cref{constr_sync:thm:prefix_before} to prefix codes.
Then, \cref{constr_sync:prop:Gamma_suffix}
is concerned with concatenation to the back.
The conditions presented here are used to deduce hardness results in later sections.

The reductions of the present section involve operations
as mentioned in \cref{prelim:rem:reductionsAC0}
and can be realized by $\AC^0$ circuits.

%\subsubsection{Concatenating to the Front}

\begin{proposition}
	\label{constrsync:prop:UanL}
	Let $L \subseteq \Sigma^*$ 
	and $U \subseteq \Gamma^*$ for some $\Gamma \subseteq \Sigma$
	be regular languages and $w \in \Sigma^*$.
	Suppose one of the following holds:
	\begin{enumerate}
	\item $L \cap \Gamma\Sigma^* = \emptyset$ with $w \in \Gamma^*$, or
	
	\item $w \notin \Gamma^*$.
	\end{enumerate}
	Then $\constrsync{wL} \leACZEROm \constrsync{UwL}$.
\end{proposition}         
\begin{proof} 
	First, assume $L \cap \Gamma\Sigma^* = \emptyset$ and $w \in \Gamma^*$.	
	Let $\semiautomat{A}{Q}{\Sigma}{\delta}$. We construct a semiautomaton $\mathcal A'$ that has a synchronizing
	word in $UwL$ if and only if $\mathcal A$ has a synchronizing word in $wL$.
	
	\medskip
	\noindent\underline{Construction:} Let $\semiautomat{A'}{Q \times \{0,1\}}{\Sigma}{\delta'}$
	where, for $i \in \{0,1\}$, $q \in Q$ and $x \in \Sigma$,
	\[
	\applytrans{\delta'}{(q, i)}{x} = 
	\begin{cases}
		(q, 0) & \text{if } i = 0 \text{ and } x \in \Gamma, \\
		(\applytrans{\delta}{q}{wx}, 1) & \text{if } i = 0 \text{ and } x \notin \Gamma, \\
		(\applytrans{\delta}{q}{x}, 1) & \text{if } i = 1.
	\end{cases}
	\]
	
	\noindent\underline{Properties:} Let $q \in Q$, $u \in U$ and $v \in L$.
	Then, as $v \notin \Gamma\Sigma^*$ and $uw \in \Gamma^*$,
	we have 
	\begin{align}
		\applytrans{\delta'}{(q, 0)}{uwv} & = \applytrans{\delta'}{(q, 0)}{v}, \label{constrsync:eqn:UL1} \\
		\applytrans{\delta'}{(q, 0)}{v}     & = (\applytrans{\delta}{q}{wv}, 1), \label{constrsync:eqn:UL3} \\
		\applytrans{\delta'}{(q, 1)}{uwv} & = (\applytrans{\delta}{q}{uwv}, 1). \label{constrsync:eqn:UL2}
	\end{align} 
	
	\noindent\underline{Verification:} Let $u \in U$ and $v \in L$.
	Then we have
	\begin{align*}
		\applytrans{\delta'}{Q \times \{0,1\}}{uwv} 
		& = \applytrans{\delta'}{Q \times \{0\}}{uwv} \cup \applytrans{\delta'}{Q \times \{1\}}{uwv} \\
		& = \applytrans{\delta'}{Q \times \{0\}}{v} \cup \applytrans{\delta'}{Q \times \{1\}}{uwv} & \tag{Eqn.~\eqref{constrsync:eqn:UL1}} \\
		& = \delta(Q, wv) \times \{1\} \cup \delta'(Q \times \{1\}, uwv) & \tag{Eqn.~\eqref{constrsync:eqn:UL3}} \\
		& = \delta(Q, wv) \times \{1\} \cup \delta(Q, uwv) \times \{1\} & \tag{Eqn.~\eqref{constrsync:eqn:UL2}} \\
		& = \delta(Q, wv)  \times \{1\}.                                   & \tag{$\delta(\delta(Q, u), wv) \subseteq  \delta(Q, wv)$}
	\end{align*}
	So $|\applytrans{\delta'}{Q \times \{0,1\}}{uwv}| = 1$
	is equivalent to $|\delta(Q,wv)| = 1$
	and $\mathcal A'$ has a synchronizing word in $UwL$
	if and only if $\mathcal A$ has a synchronizing word in $wL$.
	
    \medskip 
    
    Lastly, for the remaining case, assume $w \notin \Gamma^*$.
    Then $w = ubv$  for some $u \in \Gamma^*$, $b \notin \Gamma$, and $v \in \Sigma^*$.
    Set $L' = bvL$. We have $L' \cap \Gamma\Sigma^* = \emptyset$. 
    So, applying the previous case to $L', u$ and $U$,
    we get 
    $\constrsync{uL'} \leACZEROm \constrsync{UuL'}$
    where $uL' = wL$.
\end{proof}             

\begin{example}\label{constrsync:ex:PSPACEhard}
	By~\crefitem{constr_sync:thm:introductory_examples}{constrsync:ex:item:abcL},
	$\constrsync{c(a+b)^*}$	is \PSPACE-complete.
	By \cref{constrsync:prop:UanL}, the problems
	$\constrsync{aabbc(a+b)^*}$ and $\constrsync{a^*c(a+b)^*}$
	are \PSPACE-complete as well.
\end{example}

For reference, we list two special cases  of~\cref{constrsync:prop:UanL}.

\begin{corollary}
	\label{constrsync:cor:UanL}
	Let $L \subseteq \Sigma^*$ be regular.
	
	\begin{enumerate}
		\item \label{constrsync:cor:UanL:1} If $L \cap \Gamma\Sigma^* = \emptyset$,
		$\Gamma \subseteq \Sigma$ and $U \subseteq \Gamma^*$ is regular, then $\constrsync{L} \leACZEROm \constrsync{UL}$. 
		
		\item \label{constrsync:cor:UanL:2} $\constrsync{wL} \leACZEROm \constrsync{a^*wL}$	for $w \notin a^*$ and $a \in \Sigma$.
	\end{enumerate}
\end{corollary}                    

Next, we present a condition related to prefix codes.

\begin{theorem}
	\label{constr_sync:thm:prefix_before}
	\label{constr_sync:thm:in_font}
	%Let $U, V \subseteq \Sigma^*$ be regular and $P = U \setminus U\Sigma^+$ (those words in $U$ 
	%without a proper prefix in $U$).
	%Then $\constrsync{V} \le_m^{\log} \constrsync{PV}$. 	
	Let $L \subseteq \Sigma^*$ be regular 
	and $P \subseteq \Sigma^*$ be a prefix code that is a regular language. 
	Then 
	$
	\constrsync{L} \leACZEROm \constrsync{PL}.
	$	
\end{theorem}
\begin{proof}    	
	By~\crefitem{constr_sync:prop:aut_prefix_free}{constr_sync:prop:aut_prefix_free:item:trim} 
	there exists a trim PDFA
	$\automat{C}{R}{\Sigma}{\rho}{r_0}{\{e\}}$
	such that $P = L(\mathcal C)$
	and where the state~$e$ has no outgoing transitions.
	Suppose $\semiautomat{A}{Q}{\Sigma}{\delta}$ is an input semiautomaton. 
	
	
	% bild oder beispiel?
	\medskip 	
	\noindent\underline{Construction:} Let $\semiautomat{A'}{R \times Q}{\Sigma}{\delta'}$
	where, for $(r,q) \in R \times Q$ and $x \in \Sigma$,
	\[
	\applytrans{\delta'}{(r, q)}{x} = \begin{cases}
		(\applytrans{\rho}{r}{x}, q)   & \text{if } \applytrans{\rho}{p}{x} \text{ is defined,} \\
		(e, \applytrans{\delta}{q}{x}) & \text{if } \rho(p, x) \text{ is undefined. }
	\end{cases} 
	\]
	
	
	
	 %Property~\ref{constrsync:prefix:prop:e} readily follows from the definition of $\delta'$.
	%For the other properties, we give the proofs.  
	\medskip 	
	\noindent\underline{Properties:} Property~\ref{constrsync:prefix:defdeltastar} follows readily by induction.
	We give the proof for Property~\ref{constrsync:prefix:prop:exactsync}.
	\begin{enumerate}    	
		
		\item \label{constrsync:prefix:defdeltastar} For $r \in R$, $q \in Q$ and $u \in \Sigma^*$ we have
		\begin{equation}\label{constrsync:prefix:eqn:transitionword}
			\delta'((r, q), u)
		=
		\begin{cases}
		(e, \delta(q, xu_2)) & \text{if } \exists u_1, u_2 \in \Sigma^* \ \exists x \in \Sigma \colon u = u_1 x u_2 \land \\ 
		                    & \qquad \rho(r, u_1) \text{ is defined } \land \\ 
		                    & \qquad \rho(r, u_1 x) \text{ is undefined}, \\
		(\rho(r, u), q)     & \text{if } \rho(r, u) \text{ is defined. } 
		\end{cases}
		\end{equation}
				
\begin{comment} 				
		\item \label{constrsync:prefix:defdeltastar}
		For $r \in R$, $q \in Q$ and $u \in \Sigma^*$ we have	
		\[
		\delta'((r, q), u)
		=
		\begin{cases}
		(e, \delta(q, u)) & \text{if } r = e, \\
		(\rho(r, u), q) & \text{if } r \ne e \text{ and } \rho(r, u) \text{ is defined. }
		\end{cases}
		\]
		
		\begin{proofitem}{\ref{constrsync:prefix:defdeltastar}}
		For $r = e$ the claim follows readily from the definition of $\delta'$
		and the fact that $\rho(r, x)$ is undefined for each $x \in \Sigma$.
		Next, suppose $r \ne e$ and $\rho(r,u)$ is defined. 
		We do induction on $u$. For $u = \varepsilon$
		the claim is clear. So, suppose we have $u = vx$
		with $v \in \Sigma^*$ and $x \in \Sigma$.
		Then $\rho(r, v)$ is defined and, by the definition of $\delta'$,
		we obtain $\delta'((\rho(r, v), q), x) = (\rho(r, vx), q)$
		and the claim follows by induction.		
		\end{proofitem} 
\end{comment}
\begin{comment} 
		\item \label{constrsync:prefix:prop:e} Let $u \in \Sigma^*$ and $q \in Q$. Then $\applytrans{\delta'}{(e, q)}{u} = (e, \applytrans{\delta}{q}{u})$.
		
		\item  \label{constrsync:prefix:prop:prefixP}  Let $u \in \pref{P} \setminus P$ and $q \in Q$.
		Then $\rho(r_0, u)$ is defined and $\applytrans{\delta'}{(r_0, q)}{u} = (\rho(r_0, u), q)$
		with $\rho(r_0, u) \ne e$.
		
		\begin{proofitem}{\ref{constrsync:prefix:prop:prefixP}}
		We do induction on $u$. For $u = \varepsilon$ 
		the claim is clear. So, suppose we have $u = u'x$ for some $u' \in \Sigma^*$ and $x \in \Sigma$.
		Then, $u' \in \pref{P}\setminus P$ and, by induction hypothesis, 
		$\rho(r_0, u')$
		is defined, $\delta'((r_0, q), u') = (\rho(r_0, u'), q)$ and $\rho(r_0, u') \ne e$.    	 
		So, as $P$ is a prefix code and $u$ a proper prefix,
		$\rho(r_0, u'x)$ is defined with $\rho(r_0, u'x) \ne e$
		and, by the definition of $\delta'$,
		we find 
		\[ 
		\delta'((r_0, q), u) = \delta'((\rho(r_0, u'), q), x) = (\rho(r_0, u'x), q),
		\]
		which gives the claim by induction on the structure of $u$.
		\end{proofitem}
		
		\item \label{constrsync:prefix:prop:wordP} Let $u \in P$ and $q \in Q$. Then $\applytrans{\delta'}{(r_0, q)}{u} = (e, q)$.
		
		\begin{proofitem}{\ref{constrsync:prefix:prop:wordP}}
		Let $u = u'x$ for $u' \in \pref{P}$ and $x \in \Sigma$.
		By Property~\ref{constrsync:prefix:prop:prefixP}
		we get $\applytrans{\delta'}{(r_0, q)}{u'} = (\rho(r_0,u'), q)$
		with $\rho(r_0, u') \ne e$.
		Then, by the definition of $\delta'$, this yields
		$\applytrans{\delta'}{(r_0, q)}{u} = \applytrans{\delta'}{(\applytrans{\rho}{r_0}{u'}, q)}{x}
		= (\applytrans{\rho}{\applytrans{\rho}{r_0}{u'}}{x}, q) = (e, q)$.
		\end{proofitem}
\end{comment}
\begin{comment} 
		\item \label{constrsync:prefix:prop:undefined} Let $w \in \Sigma^*$ and $r \in R$. If $\rho(r, w)$ is undefined,
		then $\delta'((r, q), w) \in \{e\} \times Q$ for each $q \in Q$.
		
		\begin{proofitem}{\ref{constrsync:prefix:prop:undefined}}
		Let $w = w_1 x w_2$ with $w_1, w_2 \in \Sigma^*$, $x \in \Sigma$ and 
		$w_1$ of maximal length
		such that $\rho(r, w_1)$ is defined.
		By Property~\ref{constrsync:prefix:defdeltastar} 
		we have $\delta'((r, q), w_1) = (\rho(r, w_1), q)$.
		Then, by the definition of $\delta'$, as $\rho(\rho(r, w_1), x)$ is undefined,
		we find 
		\[
		\delta'((\rho(r, w_1), q), x) = (e, \delta(q, x)).
		\]
		By Property~\ref{constrsync:prefix:defdeltastar} 
		we obtain $\delta'((e, \delta(q, x)), w_2) = (e, \delta(q, xw_2)) \in \{e\} \times Q$.
		\end{proofitem}
\end{comment}
		
		\item \label{constrsync:prefix:prop:exactsync} Let $w \in \Sigma^*$ be such that $\applytrans{\rho}{R}{w} = \{e\}$.
		Then $\applytrans{\delta'}{R \times Q}{w} = \{e\} \times Q$.
		
		
		\Cref{constrsync:prefix:eqn:transitionword} 
		gives $\delta'(R \times Q, w) \subseteq \{e\} \times Q$.
		Furthermore, there exists $r \in R$ such that $\rho(r, w) = e$.
		Then, for each $q \in Q$, we have $\delta'((r, q), w) = (e,q)$
		and this gives $\{e\} \times Q = \delta'(\{r\} \times Q, w) \subseteq \delta'(R \times Q, w)$. 	    	
	\end{enumerate}   
	
	\medskip 
	\noindent\underline{Verification:} 
	By Property~\ref{constrsync:prefix:defdeltastar}
	we have $\delta'(\{r_0\} \times Q, uv) = \{e\} \times \delta(Q, v)$
	for $u \in P$ and $v \in \Sigma^*$. (note that, since~$P$ is a 
	prefix code, $uv$ can only be factored into a word from~$P$
	and an arbitrary word in one way)
	So, $|\delta'(R \times Q, uv)| = 1$
	for $u \in P$ and $v \in L$
	implies $|\delta(Q, v)| = 1$.	
	
	Conversely, suppose $|\delta(Q, v)| = 1$
	for $v \in L$. 
	By Property~\ref{constrsync:prefix:defdeltastar} and as $\rho(e,x)$ is undefined for each $x \in \Sigma$
	we have $\delta'(\{e\} \times Q, v) = \{e\} \times \delta(Q, v)$ (here, we used that~$P$ is a prefix code).
	By \crefitem{constr_sync:prop:aut_prefix_free_exact_sync}{constr_sync:prop:aut_prefix_free:item:sync} 
	there exists $w \in P$ with
	$\rho(R, w) = \{e\}$.
	Then, by Property~\ref{constrsync:prefix:prop:exactsync} 
	we have $\applytrans{\delta'}{R \times Q}{w} \subseteq \{e\} \times Q$ (note that
	we might have that $\rho(r, w)$ is undefined for some $r \in R$ 
	and that, in fact, we have equality since $\delta'(\{r_0\}\times Q), w) = \{e\} \times Q$).
	So $|\delta'(R \times Q, wv)| = 1$.
\end{proof} 

\begin{remark}
	For $L = \{c\}$ and $P = \{a,b\}^*c$
	we have that $\constrsync{L}$ 
	and $\constrsync{PL}$
	are in $\AC^0$
	and \PSPACE-complete, respectively,
	by \cref{constr_sync:thm:introductory_examples}.
	However, the complexity does not have to increase when concatenating
	a prefix code that is regular to the front of a constraint
	language as can be seen with $L = \Sigma^*$.
	Then $\constrsync{L} \eqACZEROm \constrsync{PL}$
	by \cref{constrsync:genresults:cor:left-right-ideal}.		
\end{remark}

As $\{u\}$ for $u \in \Sigma^*$ is a prefix code, we obtain the following corollary.

\begin{corollary}
	Let $L \subseteq \Sigma^*$ be regular and $u \in \Sigma^*$.
	Then $\constrsync{L} \leACZEROm \constrsync{uL}$.
\end{corollary}

%\subsubsection{Concatenating to the Back} 

Next, we state a result on concatenating
languages to the back of the constraint language.

\begin{proposition}
	\label{constr_sync:prop:Gamma_suffix}
	Let $L \subseteq \Sigma^*$ be regular and $\Gamma \subseteq \Sigma$
	such that $L \cap \Sigma^* \Gamma = \emptyset$.
	If $U \subseteq \Gamma^*$ is a regular non-empty language,
	then $\constrsync{L} \leACZEROm \constrsync{LU}$.   
\end{proposition}
\begin{proof}
	We can assume $L \ne \emptyset$. 
	Let $\semiautomat{A}{Q}{\Sigma}{\delta}$. We construct a semiautomaton $\mathcal A'$ that has a synchronizing
	word in $LU$ if and only if $\mathcal A$ has a synchronizing word in $L$.
	\FIG~\ref{fig:irred-reduction} shows the basic idea of the construction graphically.
	\begin{comment} 
		For each $q \in Q$ and $a \in \Gamma$, set $S_{q,a} = \{ \applytrans{\delta}{q}{a^i} \mid i \ge 0 \}$.
		Define  
		$\semiautomat{A'}{Q'}{\Sigma}{\delta'}$ with
		\[
		Q' = \bigcup_{q \in Q} \bigcup_{a \in \Gamma} \{q\} \times S_{q,a}.
		\]
		and, for $(q, p) \in \{q\} \times S_{q,a}$, $x \in \Sigma$
		and $p = \applytrans{\delta}{q}{a^i}$, \todo{repräsentantenunabhgängig zeigen.}
		\[
		\applytrans{\delta'}{(q, p)}{x} = \begin{cases}
			(q, \applytrans{\delta}{p}{a})     & \text{if } x = a \\
			(\applytrans{\delta}{q}{a^ix}, \applytrans{\delta}{q}{a^ix})  & \text{if } x \ne a.
		\end{cases}
		\]
		
		Let $uv$ with $u \in L$ and $v \in U$ be a synchronizing word for $\mathcal A'$.
		Note that $u$ does not end with a letter in $\Gamma$.
		So, if $|u| > 0$,\todo{genauer?}
		we have $\applytrans{\delta'}{Q'}{u}\subseteq \{ (q,q) \mid q \in Q \}$.
		For $(q,q)$ we have $\applytrans{\delta'}{(q,q)}{a^i} = (q, \applytrans{\delta}{q}{a^i})$
		for all $i \ge 0$.
		Hence, $|\applytrans{\delta'}{Q'}{u}| = |\applytrans{\delta'}{(q,q)}{v}|$
	\end{comment} 
	
	\begin{figure}[t]
		\centering
		{			
			\begin{tikzpicture}[rounded corners=10pt]						
				% Main automaton
				\node[yshift=1cm] (A) at (0,0)
				{$\mathcal A'_{|\{(q,q)\mid q \in Q\},\,\Sigma\setminus\Gamma}
					\cong
					\mathcal A_{|\Sigma\setminus\Gamma}$};
				
				\node[draw, dashed, rectangle, fit={(-5,1.4) (5,-1.4)}] {};
				
				% Left column
				\node (q) at (-3.5,-1) {$(p,p)$};				
				\node[draw,circle,inner sep=2pt] (dummyq) at (-3.5,-2.5) {};
				\node at (-3.5,-2.75) {$\vdots$};				
				\node (labSq) at (-3.5,-3.7) {$\{p\}\times Q$};
				\node[draw,dashed,rectangle,inner sep=14pt, fit={(q) (labSq)}] (Sq) {};
				
				% Right column
				\node (p) at (3.5,-1) {$(q,q)$};				
				\node[draw,circle,inner sep=2pt] (dummyp) at (3.5,-2.5) {};
				\node at (3.5,-2.75) {$\vdots$};				
				\node (labSp) at (3.5,-3.7)	{$\{q\}\times Q$};
				\node[draw,dashed,rectangle,inner sep=14pt,	fit={(p) (labSp)}] (Sp) {};
				
				% Target states
				\node[draw,circle,inner sep=2pt] (r)  at (-1.3,-0.4) {};				
				\node[draw,circle,inner sep=2pt] (r2) at (1.3, -0.4) {};
				\node[above=0.15cm of r,  xshift=-0.4cm] {$(\delta(p,ua),\delta(p,ua))$};
				\node[above=0.15cm of r2, xshift=+0.4cm] {$(\delta(q,ua),\delta(q,ua))$};
				
				% Arrows
				\path[->] (q) edge[left,pos=.7] node {$u\in\Gamma^*$} (dummyq);				
				\path[->] (p) edge[right,pos=.7] node {$u\in\Gamma^*$} (dummyp);
				
				\path[->] (dummyq) edge[bend right,below,pos=.4] node [fill=white, anchor=center] 
				{$a \in \Sigma \setminus \Gamma$} (r);
				\path[->] (dummyp) edge[bend left,below,pos=.4]  node [fill=white, anchor=center] 
				{$a \in \Sigma \setminus \Gamma$} (r2);			
			\end{tikzpicture}
					
		}
		\caption[Schematic of reduction from the proof of Proposition~\ref{constr_sync:prop:Gamma_suffix}]{Schematic of the reduction from the proof of Proposition~\ref{constr_sync:prop:Gamma_suffix}.
			For a state $q$ and word $u \in \Gamma^*$ we ``branch out'' into a copy $\{q\} \times Q$.
			For a letter from $a \in \Sigma \setminus \Gamma$,
			we return to the state $(\applytrans{\delta}{q}{ua}, \applytrans{\delta}{q}{ua})$.
%			corresponding to the state in $\mathcal A$ when reading~$ua$.
			This is indicated 
			for two states $p$ and $q$.
		}
		\label{fig:irred-reduction}
	\end{figure}
	
	
	\medskip
	\noindent\underline{Construction:} Let 
	$\semiautomat{A'}{Q'}{\Sigma}{\delta'}$
	where\footnote{Actually, setting
		\[
		Q' = \bigcup_{q \in Q} \{q\} \times S_q
		\]
		where $S_{q} = \{ \applytrans{\delta}{q}{u} \mid u \in \Gamma^* \}$
		for $q \in Q$
		would be sufficient. However, it is not clear whether the
		sets $S_q$ can be computed by an \LOGSPACE-uniform $\AC^0$ circuit family
		whereas $Q \times Q$
		can be computed by such a circuit family.} 
	$Q' = Q\times Q$
	and, for $(q, p) \in \{q\} \times Q$, $x \in \Sigma$,
	\[
	\applytrans{\delta'}{(q, p)}{x} = \begin{cases}
		(q, \applytrans{\delta}{p}{x})     & \text{if } x \in \Gamma, \\
		(\applytrans{\delta}{p}{x}, \applytrans{\delta}{p}{x})  & \text{if } x \notin \Gamma.
	\end{cases}
	\]
		
	\noindent\underline{Properties:} Next, we collect some properties of the given construction.	
	\begin{enumerate} 
		
		\item \label{constr_sync:item:word_in_Gamma} For $w \in \Gamma^*$ and $q,p \in Q$, we have
		$
		%\begin{equation}\label{constr_sync:eqn:word_in_Gamma}
		\applytrans{\delta'}{(q,p)}{w}
		= (q, \applytrans{\delta}{p}{w}).
		%\end{equation} 
		$
		
		\begin{proofitem}{\ref{constr_sync:item:word_in_Gamma}} 
		This is clear for $w = \varepsilon$.
		Next, let $w = ux$ for some $x \in \Sigma$ and $u \in \Sigma^*$.
		Inductively, we have $\applytrans{\delta'}{(q,p)}{ux}
		= \applytrans{\delta'}{(q,\delta(p, u))}{x}$.
		Since $x \in \Gamma$,
		by the definition of the transition function
		we find $\applytrans{\delta'}{(q,\delta(p, u))}{x} = (q, \delta(p,w))$.			
		\end{proofitem}		
		
		\item \label{constr_sync:item:Gamma_suffix} If $w \in \Sigma^* (\Sigma \setminus \Gamma)$ and $q,p \in Q$, then
		$
		%\begin{equation}\label{constr_sync:eqn:Gamma_suffix}
		\applytrans{\delta'}{(q,p)}{w} = (\applytrans{\delta}{p}{w}, \applytrans{\delta}{p}{w}).
		%
		%\end{equation}
		$
				
		\begin{proofitem}{\ref{constr_sync:item:Gamma_suffix}} 
		By assumption, we can write $w = w_0 u_1 w_1 \cdots u_n w_n$ for some $n \ge 0$
		with $w_0, w_1, \ldots, w_n \in (\Sigma \setminus \Gamma)^+$
		and $u_1, \ldots, u_n \in \Gamma^*$. For $n = 0$,
		the claim follows by an easy induction from the defining equation of~$\delta'$.
		Next, suppose $n \ge 1$. Since $|w_{n-1}| \ge 1$,
		inductively, we obtain
		\[
		\delta'((q,p), w') 
		= (\delta(p, w'), \delta(p, w'))
		\]
		for $w' = w_0 u_1 w_1 \cdots w_{n-1}$. By Property~\ref{constr_sync:item:word_in_Gamma}
		we have
		\[
		\delta'((\delta(p, w'), \delta(p, w')), u_n) = (\delta(p, w'), \delta(p, w'u_n)).
		\]
		Then applying the claim for $n = 0$ again, we find 
		\[ 
		\delta'((\delta(p, w'), \delta(p, w'u_n)), w_n) = (\delta(p, w'u_n w_n), \delta(p, w'u_n w_n)).
		\]
		Thus we have
		$
		\applytrans{\delta'}{(q,p)}{w} = (\applytrans{\delta}{p}{w}, \applytrans{\delta}{p}{w}).
		$		
		\end{proofitem}		
	
		\item \label{constr_sync:item:QxQ}  For $v \in L \cap \Sigma^+$, 
		we have  $\delta'(Q', v) = \{(q,q) \mid q \in \delta(Q, v) \}$.
		
		\begin{proofitem}{\ref{constr_sync:item:QxQ}}
		By assumption on~$L$,
		we have $L \cap \Sigma^+ \subseteq \Sigma^*( \Sigma \setminus \Gamma )$
		and so $v \in \Sigma^* (\Sigma \setminus \Gamma)$.
		Thus, the claim follows from Property~\ref{constr_sync:item:Gamma_suffix}.		
		\end{proofitem}
	\end{enumerate}	
	
	\medskip 
	\noindent\underline{Verification:} We can assume $|Q| \ge 2$.
	First, suppose $|\delta(Q, v) | = 1$ for some $v \in L$.
	Then, by Property~\ref{constr_sync:item:QxQ}, $|\delta'(Q \times Q, v)| = 1$.
	So, $|\delta'(Q', vu)| = 1$ for $u \in U$.	
	
	Conversely, assume $|\delta'(Q', vu)| = 1$
	for $v \in L$ and $u \in U$. %, where we can assume $|v| > 0$.
%	By Property~\ref{constr_sync:item:Gamma_suffix}, as $v \notin \Sigma^*\Gamma$,
%	we have $\delta'(Q\times Q, v) = \{(q,q) \mid q \in \delta(Q, v) \}$.
	By Property~\ref{constr_sync:item:word_in_Gamma}, for $p,q \in Q$ and $u \in \Gamma^*$,
	we have $\delta'((q,q), u) = (q, \delta(q, u))$ and $\delta'((p,p), u) = (p, \delta(p, u))$.
	So,
	\[
	\delta'((q,q), u) = \delta'((p,p), u)
	\Leftrightarrow 
	q = p.
	\]
	Hence, with Property~\ref{constr_sync:item:QxQ}, $|\delta'(Q', vu)| = 1$
	implies $|\delta'(Q', v)| = 1$,
	which, again by Property~\ref{constr_sync:item:QxQ}, then implies $|\delta(Q, v)| = 1$.	
%	
%	
%	
%	
%	
%	
%	\bigskip 
%	
%	Let $v \in L$ and $u \in U$. We can assume $|v| > 0$.
%	By Property~\ref{constr_sync:item:Gamma_suffix}, as $v \notin \Sigma^*\Gamma$,
%	we have $\delta'(Q\times Q, v) = \{(q,q) \mid q \in \delta(Q, v) \}$.
%	So,  $|\delta(Q, v) | = 1$,
%	implies $|\delta'(Q \times Q, vu)| = 1$.
%	
%	%By Property~\ref{} for $v \in \Gamma^*$ 
%	%we have $\delta'(S \times T, v) = \bigcup_{q \in S} \{q\} \times \delta(T, v)$.
%	
%	Furthermore, by Property~\ref{constr_sync:item:word_in_Gamma}, as $u \in \Gamma^*$, we have $(q, \delta(q, u)) = \delta'((q,q), u)$ and $(p, \delta(p, u)) = \delta'((p,p), u)$,
%	and so $\delta'((q,q), u) = \delta'((p,p), u) = (p, \delta(u, p))$
%	if and only if $q = p$ for $q,p \in Q$. This yields
%	that $|\delta'(Q \times Q, vu)| = 1$
%	implies $|\delta'(Q \times Q, v)| = 1$,
%	which then implies $|\delta(Q, v)| = 1$.
\end{proof} 

Lastly, we give a condition for preservation of the complexity w.r.t.~concatenation.

\begin{proposition} \label{constrsync:genresults:prop:concat}
	Let $L_1, L_2, L_3 \subseteq \Sigma^*$ be regular languages.
	If $L_1 L_2 L_3 \subseteq \fact{L_2}$, then $
	\constrsync{L_1 L_2 L_3} \eqACZEROm \constrsync{L_2}.
	$
\end{proposition}
\begin{proof}
	In general, it is $\fact{L_2} 
	\subseteq 
	\fact{L_1 L_2 L_3}$.
	By assumption, $\fact{L_1 L_2 L_3}
	\subseteq \fact{L_2}$.
	So, the claim follows from~\cref{gen_results:cor:factors_etc_complexity}.
\end{proof}

\begin{corollary}
 	For $u,v \in \Sigma^*$ and $L \subseteq \Sigma^*$ with $uLv \subseteq \fact{L}$
 	we have $\constrsync{uLv} \eqACZEROm \constrsync{L}$.
\end{corollary}